\subsubsection{Acknowledge (ACK)}
\label{chap:Back_I2C_Protocol_ACK}

After every transmitted byte, there is an acknowledge phase. The ACK bit is 
used by the receiver, whereby the receiver can signal the transmitter that the byte 
was received (not the byte, rather the byte is all byte transfers, i.e. signal the 
transmitter that the byte transfer was completed). The master generates all clock 
pulses, including the ninth clock cycle, which is used for ACK/NACK. During this 
ninth clock, the transmitter releases SDA so the receiver can drive it. 

An ACK is signaled when the receiver pulls SDA LOW and holds it while SCL 
is HIGH for that clock cycle. A NACK is signaled when SDA stays HIGH during 
the ninth clock cycle. 

After a NACK, the controller usually either does a STOP to end the transfer 
or a repeated START to initiate a new transaction. NACKs can occur for various 
reasons, such as the absence of a device at the transmitted address, the device 
that was addressed is busy and cannot respond, the receiver is unable to process 
the command or data that was sent, the receiver has a busy status meaning it 
cannot accept additional data, or a controller-receiver pair uses NACK to signal a 
termination of the transfer to the target transmitter.